\noindent\textbf{Tämän kirjan omistaa:}
\\
\\
\\
$\cdots\cdots\cdots\cdots\cdots\cdots\cdots\cdots\cdots\cdots$
\\


\newpage

\begin{figure}[h!]
\centering
\includegraphics[scale=0.4]{graphics/logo.png}
\end{figure}

\cleardoublepage




\textbf{Julkaisija:}

Tekis
\\


\textbf{Laulukirjavastaavat:}

Julius Uusinarkaus

Maximillian Remming
\\

\textbf{Taittosofta:}

PS

Max
\\

\textbf{Lauluja keräsi:}

Felix, Julius, Max, jne, jne 
\\

\textbf{Suurkiitokset:}

Sheets, Python, \LaTeX
\\

\textbf{Painos:}

5/2017 - 520 kpl

\newpage

\textbf{Esipuhe}
\\

Matemaatikot latoivat laulukirjan käsin latexilla, perinteitä kunnioittaen.
Valtiotieteilijät toisaalta pitkillä kokouksilla. Itse kunnioitimme
tietojenkäsittelijöiden perinteitä. Laulukirjamme on ladottu excelistä
python-scriptillä latexiin. Sopivalla määrällä meemitaikaa. Niin, ja tätä on työstetty 6 vuotta.

Laulukirja on paljon enemmän kuin laulukirja. Kirja on tietenkin hyvä työkalu
sitseillä, mutta sen aito tarkoitus on jäädä raportiksi menneestä. Laulukirjaan
kerääntyy satoja viestejä, rivouksia ja muita tervehdyksiä, nykyisiltä ja
menneiltä ystäviltä. Laulukirja kuluu, mutta paranee iän myötä. Teimme tämän
laulukirjan sellaiseksi, että siitä on sinulle vielä iloa pitkään. Siitä tulee
jokaisen sitsin jälkeen enemmän omannäköisesi. Voit lisätä puuttuvat lempilaulusi DLC-sivuille.

We wanted to make the book accessible, and actually useful, even if you don’t
speak finnish. We put together as many “fun to drink together, or while drunk”-songs
in english as we could. We tried to make this the ultimate combination of a
sophisticated traditional song book, and meme magic. I think we did well.
\\
\\
Hyvää vappua!

\vfill
\textit{Julius ja Max}

\cleardoublepage
